\documentclass[11pt,aspectratio=169]{beamer}

% ============================================================
%  MODULE: Descriptive Statistics & Frequentist Statistics
%  Detailed, plain-style study notes (no colour boxes)
%  Includes expanded content + full speaker notes on every slide
% ============================================================

\usetheme{default}
\usecolortheme{dove}          % grayscale, plain, no colour highlight boxes
\setbeamertemplate{blocks}[default]   % plain framed blocks, no coloured fill
\setbeamercolor{block title}{bg=white,fg=black}
\setbeamercolor{block body}{bg=white,fg=black}
\setbeamercolor{frametitle}{bg=white,fg=black}
\setbeamercolor{title}{bg=white,fg=black}
\setbeamercolor{structure}{fg=black}
\setbeamertemplate{itemize items}[default]
\setbeamertemplate{navigation symbols}{}
\setbeamertemplate{footline}[frame number]

\usepackage[utf8]{inputenc}
\usepackage{amsmath,amssymb,amsthm}
\usepackage{mathtools}
\usepackage{bm}
\usepackage{booktabs}

\title[Descriptive \& Frequentist Statistics]{Descriptive Statistics and Frequentist Statistics}
\subtitle{Module 3}
\author{Aishwarya J S}


\newcommand{\Var}{\operatorname{Var}}
\newcommand{\Cov}{\operatorname{Cov}}
\newcommand{\Bias}{\operatorname{Bias}}
\newcommand{\MSE}{\operatorname{MSE}}
\newcommand{\E}{\operatorname{\mathbb{E}}}
\newcommand{\Xbar}{\bar{X}}

% Use plain \paragraph-style headers instead of coloured blocks
\newcommand{\Def}[1]{\textbf{Definition.} #1}
\newcommand{\Note}[1]{\textbf{Note.} #1}
\newcommand{\Res}[1]{\textbf{Result.} #1}
\newcommand{\Ex}[1]{\textbf{Example.} #1}

\begin{document}

\begin{frame}
\titlepage
\note{Title slide. This deck is the full module: Descriptive Statistics
(histogram, mean/variance, order statistics, covariance, covariance matrix)
and Frequentist Statistics (sampling, MSE, consistency, confidence
intervals, parametric and non-parametric estimation). Every slide's speaker
note repeats and extends the on-slide content so this can be read as a
stand-alone study document, not just a slide skeleton.}
\end{frame}

\begin{frame}{Module Outline}
\tableofcontents
\note{Two main sections. Descriptive statistics is about summarizing data
you already have, with no reference to an underlying random model beyond
what is needed to justify formulas. Frequentist statistics is about using a
sample to infer properties of an assumed underlying population/model, under
the interpretation that probability is long-run frequency and the unknown
parameter is a fixed constant, not a random variable.}
\end{frame}

% ============================================================
\section{Descriptive Statistics}
% ============================================================

\begin{frame}{What is Descriptive Statistics?}
\Def{Descriptive statistics summarizes and organizes the features of a data
set using numerical measures, tables, and graphs, without drawing
conclusions that go beyond the observed data.}

\begin{itemize}
    \item Goal: describe \textbf{center}, \textbf{spread}, \textbf{shape},
          and \textbf{relationships}.
    \item Graphical summaries: histogram, boxplot, scatterplot, stem-and-leaf
          plot, bar chart, pie chart.
    \item Numerical summaries: mean, median, mode, variance, standard
          deviation, order statistics, covariance, correlation.
    \item Descriptive statistics makes \emph{no probabilistic claim} about a
          population; it only reports what is present in the sample.
          Contrast this with inferential (frequentist) statistics, where we
          use the sample to say something about a population we have not
          fully observed.
\end{itemize}
\note{Key exam point: descriptive vs.\ inferential distinction. Descriptive
statistics can be computed on a finite, fixed data set with no assumptions
about randomness at all — e.g.\ you can describe the marks of exactly the
30 students in a class without any probability model. Frequentist/inferential
statistics requires you to view the data as one realization of a random
sample from a larger population, and it is this second view that lets you
attach a sampling distribution, standard error, or confidence interval to a
statistic. Also note the four descriptive dimensions (center, spread,
shape, relationship) — most exam questions map directly onto one of these
four.}
\end{frame}

% ---------------- HISTOGRAM ----------------
\subsection{Histogram}

\begin{frame}{Histogram: Definition and Construction}
\Def{A histogram is a graphical estimate of the distribution of a
numerical variable, built by (1) partitioning the range of the data into
contiguous bins, (2) counting the number of observations in each bin, and
(3) drawing bars whose \emph{area} is proportional to that count.}

For equal-width bins of width $h$, the \textbf{density-scaled} height is
\[
      \hat{f}(x) = \frac{n_i}{n h}, \qquad x \in [a_{i-1}, a_i)
\]
so that the total area under the histogram is exactly $1$, mimicking a
probability density function.

\Note{A frequency histogram (raw counts on the $y$-axis) and a density
histogram (as above) have the same shape; only the vertical scale differs.
Density scaling is what allows the histogram to be compared directly with a
theoretical pdf overlay.}
\note{Walk through the three-step construction slowly: choose bins, count,
draw. Emphasize that "area proportional to frequency" is the defining
property, not "height proportional to frequency" — this only coincides
with height when all bins have equal width, which is the usual case in
practice but not a requirement. If bins have unequal widths, you MUST
scale by width or the picture is visually misleading (larger bins would
look artificially tall). This is a classic exam trap.}
\end{frame}

\begin{frame}{Histogram: Choice of Bin Width}
Bin choice is a bias--variance trade-off:
\begin{itemize}
    \item Too few bins (wide $h$) $\Rightarrow$ \textbf{oversmoothed}, hides
          real structure (multi-modality, skew) — high bias.
    \item Too many bins (narrow $h$) $\Rightarrow$ \textbf{noisy/spiky} —
          high variance, dominated by sampling noise.
\end{itemize}

Common automatic rules ($n$ = sample size, $s$ = sample std.\ dev., IQR =
interquartile range):
\[
    \text{Sturges: } k = \lceil \log_2 n \rceil + 1
    \qquad\qquad
    \text{Scott: } h = \frac{3.49\,s}{n^{1/3}}
\]
\[
    \text{Freedman--Diaconis: } h = \frac{2\,\text{IQR}}{n^{1/3}}
\]

\Note{Freedman--Diaconis is the most robust of the three because IQR is
insensitive to outliers, whereas Scott's rule (which uses $s$) can be
badly distorted by even one extreme observation.}
\note{Sturges' rule targets a specific number of bins $k$; Scott's and
Freedman--Diaconis target a bin width $h$ directly, and you convert to $k$
via $k = \lceil \text{range}/h \rceil$. Sturges assumes roughly normal data
and under-bins for large, non-normal, or heavy-tailed data sets — this is
a common critique to mention in an exam answer. All three rules scale bin
width down as $n^{-1/3}$: more data supports finer bins, but only at a
slow cube-root rate.}
\end{frame}

\begin{frame}{Histogram as a Non-Parametric Density Estimator}
\begin{itemize}
  \item The histogram is the simplest \textbf{non-parametric density
        estimator}: it estimates the unknown pdf $f(x)$ using only the data,
        with no assumed functional form for $f$.
  \item For a fixed bin $[a_{i-1},a_i)$ with true probability
        $p_i = \int_{a_{i-1}}^{a_i} f(x)\,dx$, the count $n_i$ is Binomial$(n,p_i)$, so
        $\hat f(x) = n_i/(nh)$ is unbiased for $p_i/h$ (an average of $f$
        over the bin, not exactly $f(x)$ at a point).
  \item \textbf{Consistency} of the histogram as an estimator of $f$
        requires both $h \to 0$ \emph{and} $nh \to \infty$ as
        $n \to \infty$: bins must shrink (to reduce bias) but still contain
        enough points on average (to reduce variance).
\end{itemize}
\note{This slide connects descriptive statistics to the later
non-parametric estimation section. Emphasize the two competing limits:
$h\to 0$ alone is not enough, because as bins shrink, the expected count
per bin $np_i \approx nhf(x)$ would shrink too and the estimate would
become pure noise; you additionally need $nh\to\infty$ so that the number
of points per bin still grows. Kernel density estimation, covered later, is
the smooth generalization of exactly this idea — it replaces the rigid
box-shaped bin membership indicator with a smooth kernel.}
\end{frame}

% ---------------- SAMPLE MEAN AND VARIANCE ----------------
\subsection{Sample Mean and Variance}

\begin{frame}{Sample Mean}
Let $X_1,\dots,X_n$ be a random sample from a population with mean $\mu$
and variance $\sigma^2$.
\[
    \Xbar = \frac{1}{n}\sum_{i=1}^{n} X_i
\]

Properties:
\begin{itemize}
    \item \textbf{Unbiased}: $\E[\Xbar] = \mu$.
    \item \textbf{Variance}: $\Var(\Xbar) = \sigma^2/n$ — shrinks as $n$
          grows, so $\Xbar$ becomes more precise with more data.
    \item \textbf{Law of Large Numbers (LLN)}: $\Xbar \xrightarrow{P} \mu$
          as $n \to \infty$.
    \item \textbf{Central Limit Theorem (CLT)}:
          $\sqrt{n}(\Xbar-\mu) \xrightarrow{d} N(0,\sigma^2)$, so for large
          $n$, $\Xbar \approx N(\mu, \sigma^2/n)$ \emph{regardless of the
          shape of the original population distribution}.
\end{itemize}
\note{The CLT is the single most important theorem underlying the entire
frequentist section: it is why we can use normal- or $t$-based confidence
intervals even when the underlying population is not normal, as long as
$n$ is reasonably large. Contrast: if the population itself IS normal, then
$\Xbar$ is exactly (not just approximately) normal for every $n$, even
$n=1$. Also mention: variance of $\Xbar$ decreasing like $1/n$ means
standard error (its square root) decreases like $1/\sqrt n$ — to halve the
standard error you need four times the data, a very commonly tested fact.}
\end{frame}

\begin{frame}{Sample Variance}
\[
    S^2 = \frac{1}{n-1}\sum_{i=1}^n (X_i - \Xbar)^2
    = \frac{1}{n-1}\left(\sum_{i=1}^n X_i^2 - n \Xbar^2\right)
\]

\begin{itemize}
    \item Dividing by $n-1$ (\textbf{Bessel's correction}) makes $S^2$
          \textbf{unbiased}: $\E[S^2] = \sigma^2$.
    \item The ``plug-in''/MLE-style version divides by $n$ instead:
    \[
        \hat\sigma^2 = \frac{1}{n}\sum_{i=1}^n (X_i-\Xbar)^2 = \frac{n-1}{n}S^2,
        \qquad \E[\hat\sigma^2] = \frac{n-1}{n}\sigma^2 \ \ (\text{biased low})
    \]
    \item $n-1$ = \textbf{degrees of freedom} remaining after estimating
          $\mu$ with $\Xbar$.
\end{itemize}

\textbf{Why $n-1$?} The deviations satisfy $\sum_i (X_i-\Xbar)=0$, so only
$n-1$ of them are free to vary once $\Xbar$ is fixed; using $\Xbar$ (which
is, on average, closer to each $X_i$ than the unknown true $\mu$ is) makes
$\sum(X_i-\Xbar)^2$ systematically smaller than $\sum(X_i-\mu)^2$, and
dividing by the smaller number $n-1$ exactly compensates.
\note{Be ready to reproduce the short algebraic proof that
E[sum (Xi - Xbar)^2] = (n-1) sigma^2: expand
sum(Xi-mu)^2 = sum(Xi - Xbar + Xbar - mu)^2 and use that
sum(Xi-Xbar)(Xbar-mu) cross term vanishes appropriately, and
E[(Xbar-mu)^2] = sigma^2/n, giving sum term = n sigma^2 - sigma^2 = (n-1)
sigma^2. Also mention standard deviation S = sqrt(S^2) is NOT an unbiased
estimator of sigma even though S^2 is unbiased for sigma^2 (square root is
a nonlinear/concave transform, so by Jensen's inequality E[S] < sigma in
general) — a favourite trick question.}
\end{frame}

\begin{frame}{Other Measures of Center and Spread}
\begin{itemize}
    \item \textbf{Median}: robust measure of center, unaffected by extreme
          outliers (unlike the mean).
    \item \textbf{Mode}: most frequent value; useful for categorical or
          multi-modal data.
    \item \textbf{Range}: $X_{(n)}-X_{(1)}$; highly sensitive to outliers.
    \item \textbf{Standard deviation}: $S=\sqrt{S^2}$, same units as the
          data (variance is in squared units).
    \item \textbf{Coefficient of variation}: $CV = S/\Xbar$, a unit-free
          measure of relative spread, useful for comparing variability
          across data sets with different units or scales.
    \item \textbf{Skewness} and \textbf{kurtosis}: standardized third and
          fourth central moments, describing asymmetry and tail-heaviness
          of the distribution's shape.
\end{itemize}
\note{This is additional depth beyond the syllabus wording but commonly
expected in a full descriptive-statistics module: mean/median/mode covers
"center", range/variance/sd/CV covers "spread", skewness/kurtosis covers
"shape" — completing the earlier four-way breakdown (center, spread,
shape, relationship). Mean is pulled toward outliers/skew, median is not;
this is why income data (right-skewed) is usually summarized with the
median rather than the mean. Skewness and kurtosis are only introduced
here as a bullet — the following subsection develops both in full, since
they are frequently under-covered relative to how often they are tested.}
\end{frame}

% ---------------- SKEWNESS AND KURTOSIS ----------------
\subsection{Skewness and Kurtosis}

\begin{frame}{Skewness: Definition}
\Def{Skewness measures the \emph{asymmetry} of a probability distribution
about its mean. It is the standardized third central moment.}

\textbf{Population skewness:}
\[
    \gamma_1 = \frac{\E[(X-\mu)^3]}{\sigma^3} = \frac{\mu_3}{\sigma^3}
\]

\begin{itemize}
    \item $\mu_3 = \E[(X-\mu)^3]$ is the third central moment.
    \item Dividing by $\sigma^3$ makes $\gamma_1$ \textbf{unit-free} and
          \textbf{scale-invariant} — the same idea used to turn covariance
          into a correlation coefficient.
    \item A symmetric distribution (e.g.\ Normal) has $\gamma_1=0$; cubing
          preserves sign, so $\gamma_1$ captures \emph{which side} has the
          longer tail, unlike variance (which squares deviations and loses
          all sign information).
\end{itemize}
\note{Contrast the exponents directly: variance uses the SECOND power
(always non-negative, so it can only measure spread, never direction of
asymmetry), while skewness uses the THIRD power, which keeps its sign
(a large positive deviation cubed stays positive, a large negative
deviation cubed stays negative). This is exactly why skewness, unlike
variance, can tell you whether the long tail sits to the right or the
left of the mean. Also point out this is the natural next member of the
"standardized $k$-th central moment" family that variance/SD already
belong to — kurtosis (covered a few slides later) is the $k=4$ member.}
\end{frame}

\begin{frame}{Sample Skewness: Formulas}
Given a sample $X_1,\dots,X_n$ with mean $\Xbar$:

\textbf{Method-of-moments (plug-in) sample skewness:}
\[
    g_1 = \frac{m_3}{m_2^{3/2}}, \qquad
    m_k = \frac{1}{n}\sum_{i=1}^n (X_i-\Xbar)^k
\]

\textbf{Adjusted Fisher--Pearson coefficient} (bias-corrected; used by
default in Excel, SPSS):
\[
    G_1 = \frac{\sqrt{n(n-1)}}{n-2}\, g_1
\]

\begin{itemize}
    \item $g_1$ is the ``raw'' plug-in estimator — analogous to
          $\hat\sigma^2$ (divide by $n$), it is slightly biased in finite
          samples.
    \item $G_1$ applies a correction analogous to Bessel's correction for
          variance, and is preferred for small/moderate $n$.
    \item As $n\to\infty$, $G_1 \to g_1$, and both are \textbf{consistent}
          estimators of $\gamma_1$.
\end{itemize}
\note{Draw the direct parallel to the sample-variance slide: there, the
plug-in $\hat\sigma^2$ (divide by $n$) is biased and the corrected $S^2$
(divide by $n-1$) is unbiased; here, $g_1$ is the plug-in analogue and
$G_1$ is the corrected analogue, though the correction factor for the
third moment is more involved than a simple $n/(n-1)$ ratio because
skewness is a ratio of moments, not a single moment. Students should be
able to identify $m_2$ here as essentially the population-style
($n$-divisor) variance, i.e.\ $m_2 = \hat\sigma^2$ from the sample-variance
slide, reused directly in the skewness formula.}
\end{frame}

\begin{frame}{Types of Skewness}
\begin{itemize}
    \item \textbf{Positive (right) skew:} $\gamma_1>0$. Long tail stretches
          to the right; a few large values pull the mean upward.
    \[
        \text{Mean} > \text{Median} > \text{Mode}
    \]
    Example: income distributions, house prices, waiting times.

    \item \textbf{Negative (left) skew:} $\gamma_1<0$. Long tail stretches
          to the left.
    \[
        \text{Mean} < \text{Median} < \text{Mode}
    \]
    Example: age at retirement, exam scores on an easy test.

    \item \textbf{Symmetric:} $\gamma_1=0$. Mean $\approx$ Median
          $\approx$ Mode (holds exactly for unimodal symmetric
          distributions like the Normal).
\end{itemize}

\Note{The mean--median--mode ordering above is a common heuristic, not a
theorem — it can fail for some multimodal or heavy-tailed distributions.}
\note{This slide directly connects back to the "Other Measures of Center
and Spread" slide's remark that the median is preferred over the mean for
skewed data: the ordering here explains WHY — in right-skewed data the
mean is pulled furthest toward the long tail, the mode stays anchored at
the most common (peak) value, and the median sits in between. Ask
students to sketch a right-skewed density and mark all three measures on
it as a self-check; the visual makes the ordering intuitive rather than
something to memorize by rote.}
\end{frame}

\begin{frame}{Pearson's Skewness Coefficients \& Interpreting Magnitude}
Quick, robust approximations that avoid computing third moments directly:

\textbf{Pearson's first coefficient (mode-based):}
\[
    Sk_1 = \frac{\Xbar - \text{Mode}}{S}
\]

\textbf{Pearson's second coefficient} (median-based, more commonly used
since the mode is often ill-defined for continuous data):
\[
    Sk_2 = \frac{3(\Xbar - \text{Median})}{S}
\]

\textbf{Rule-of-thumb magnitude interpretation} (for $\gamma_1$ or $g_1$):
\begin{itemize}
    \item $|\text{skewness}| < 0.5$: approximately symmetric
    \item $0.5 \le |\text{skewness}| < 1$: moderately skewed
    \item $|\text{skewness}| \ge 1$: highly skewed
\end{itemize}
\note{Pearson's coefficients are historically important because they let
you get a quick sense of skew direction and rough magnitude using only
mean, median/mode, and standard deviation — no need to compute a third
moment by hand. The factor of 3 in Pearson's second coefficient comes from
an approximate relationship, valid for moderately skewed unimodal
distributions, that (Mean - Mode) $\approx$ 3(Mean - Median); this is
where $Sk_1 \approx Sk_2$ under that approximation. The magnitude
thresholds are guidance from common statistical software/textbooks, not a
universal hard rule — flag this if a strict cutoff is asked for in an
exam, since different sources quote slightly different boundaries.}
\end{frame}

\begin{frame}{Kurtosis: Definition and Types}
\Def{Kurtosis is the standardized \emph{fourth} central moment; it
describes the \emph{tail-heaviness / peakedness} of a distribution
relative to the Normal.}
\[
    \text{Kurtosis} = \frac{\E[(X-\mu)^4]}{\sigma^4}, \qquad
    \text{Excess Kurtosis} = \text{Kurtosis} - 3
\]

\begin{itemize}
    \item \textbf{Normal distribution}: skewness $=0$, excess kurtosis
          $=0$ (\textbf{mesokurtic}) — the ``3'' is subtracted precisely so
          the Normal serves as the zero baseline.
    \item \textbf{Leptokurtic} (excess kurtosis $>0$): heavier tails, more
          outlier-prone than Normal (e.g.\ Student's $t$).
    \item \textbf{Platykurtic} (excess kurtosis $<0$): lighter tails,
          flatter peak (e.g.\ Uniform).
\end{itemize}

\Note{Skewness answers ``which side is longer?''; kurtosis answers ``how
heavy are the tails / how outlier-prone is the data?'' Together they
complete the ``shape'' dimension of descriptive statistics, alongside
center (mean/median/mode) and spread (variance/SD/IQR). This matters
practically in finance, anomaly detection, and checking the Normality
assumptions behind many of the frequentist procedures (e.g.\ $t$- and
$z$-intervals) introduced later in this module.}
\note{Emphasize the "-3" convention carefully, since it is a frequent
source of confusion: RAW kurtosis for a Normal distribution is exactly 3,
not 0; EXCESS kurtosis subtracts 3 so that Normal becomes the natural
zero-reference point, analogous to how excess return in finance is
measured relative to a risk-free baseline. Different software/textbooks
sometimes report raw kurtosis and sometimes excess kurtosis without
saying so explicitly — always check which convention a given number uses
before interpreting it. Also worth flagging: just as skewness has a
plug-in ($g_1$) and bias-corrected ($G_1$) sample version, kurtosis has an
analogous pair of sample estimators, following exactly the same
plug-in-vs-corrected pattern already seen for variance and skewness.}
\end{frame}

% ---------------- ORDER STATISTICS ----------------
\subsection{Order Statistics}

\begin{frame}{Order Statistics: Definition}
Arrange the sample $X_1,\dots,X_n$ in increasing order:
\[
    X_{(1)} \le X_{(2)} \le \cdots \le X_{(n)}
\]
$X_{(k)}$ is the $k$-th \textbf{order statistic}.

\begin{itemize}
    \item $X_{(1)}=\min$, $X_{(n)}=\max$, \ \ Range $R = X_{(n)}-X_{(1)}$.
    \item \textbf{Median}:
    $
        X_{((n+1)/2)}\ (n \text{ odd}), \quad
        \tfrac12\left(X_{(n/2)}+X_{(n/2+1)}\right)\ (n \text{ even})
    $
    \item Order statistics are \textbf{not independent} of each other
          (even if the original $X_i$ are i.i.d.), because knowing
          $X_{(1)}$ constrains what $X_{(2)},\dots$ can be.
\end{itemize}
\note{Order statistics are the foundation of robust statistics — measures
built from them (median, IQR, trimmed mean) resist distortion by a small
number of extreme or erroneous observations, unlike the mean and variance
which use every observation with equal weight and hence can be dragged
arbitrarily far by a single outlier. This robustness/efficiency trade-off
(order-statistic-based measures vs.\ moment-based measures) is a recurring
theme across the whole module.}
\end{frame}

\begin{frame}{Quantiles, Quartiles, IQR, Boxplot}
\begin{itemize}
    \item The \textbf{$p$-th sample quantile} is (approximately) the value
          below which a fraction $p$ of the data falls; computed by
          interpolating between order statistics near rank $p(n-1)+1$
          (conventions vary slightly across software).
    \item \textbf{Quartiles}: $Q_1$ (25th percentile), $Q_2=$ median (50th),
          $Q_3$ (75th percentile).
    \item \textbf{Interquartile range}: $\text{IQR}=Q_3-Q_1$ — spread of the
          middle 50\% of data, robust to outliers.
    \item \textbf{Five-number summary}: $X_{(1)}, Q_1, \text{Median}, Q_3,
          X_{(n)}$ — basis of the boxplot.
    \item \textbf{Outlier rule of thumb}: flag points outside
          $[Q_1 - 1.5\,\text{IQR},\ Q_3 + 1.5\,\text{IQR}]$ (Tukey's fences).
\end{itemize}
\note{Tukey's 1.5 x IQR fence rule is worth memorizing explicitly, it is a
frequently tested numeric detail. Boxplots visually encode the five-number
summary: box spans Q1 to Q3 with a line at the median, whiskers extend to
the most extreme non-outlier points, and points beyond the fences are
plotted individually as outliers.}
\end{frame}

\begin{frame}{Sampling Distribution of Order Statistics}
For a continuous r.v.\ with pdf $f$ and cdf $F$, the density of the $k$-th
order statistic (out of $n$ i.i.d.\ draws) is
\[
    f_{X_{(k)}}(x) = \frac{n!}{(k-1)!\,(n-k)!}\, [F(x)]^{k-1}\,[1-F(x)]^{n-k}\, f(x)
\]
Special cases:
\[
    f_{X_{(1)}}(x) = n[1-F(x)]^{n-1}f(x), \qquad
    f_{X_{(n)}}(x) = n[F(x)]^{n-1}f(x)
\]
\Ex{For $X_i \sim \text{Uniform}(0,1)$, $F(x)=x$, so
$f_{X_{(n)}}(x) = n x^{n-1}$ on $(0,1)$, giving
$\E[X_{(n)}] = n/(n+1) \to 1$ as $n\to\infty$: the sample maximum
approaches the true upper endpoint.}
\note{Derive the formula intuitively: for X_(k) to land near x, you need
exactly k-1 of the n observations below x, one observation at x, and n-k
observations above x — the multinomial coefficient n!/((k-1)!1!(n-k)!)
counts the ways to choose which points play which role, F(x)^{k-1} is the
probability the k-1 chosen points are below x, [1-F(x)]^{n-k} is the
probability the rest are above, and f(x)dx is the density contribution of
the single point at x. This distributional result underlies extreme value
theory and reliability analysis (e.g., time to first failure among n
components is min of n lifetimes).}
\end{frame}

% ---------------- SAMPLE COVARIANCE ----------------
\subsection{Sample Covariance}

\begin{frame}{Sample Covariance}
For paired data $(X_1,Y_1),\dots,(X_n,Y_n)$:
\[
    S_{XY} = \frac{1}{n-1}\sum_{i=1}^{n} (X_i-\Xbar)(Y_i-\bar Y)
    = \frac{1}{n-1}\left(\sum_i X_iY_i - n\Xbar\bar Y\right)
\]

\begin{itemize}
    \item $S_{XY}>0$: $X$ and $Y$ tend to increase/decrease together.
    \item $S_{XY}<0$: $X$ and $Y$ tend to move in opposite directions.
    \item $S_{XY}\approx 0$: little or no \emph{linear} association.
    \item Magnitude depends on the units of $X$ and $Y$, so raw covariance
          values are \textbf{not comparable} across different variable
          pairs — this motivates the correlation coefficient.
\end{itemize}
\note{Emphasize covariance only detects LINEAR association: two variables
can be strongly, deterministically related in a nonlinear way (e.g.
Y = X^2 for X symmetric around 0) and still have covariance exactly zero.
This is a very common conceptual pitfall to highlight. Also note S_XY is
symmetric: S_XY = S_YX.}
\end{frame}

\begin{frame}{Sample Correlation Coefficient}
\[
    r = \frac{S_{XY}}{S_X S_Y} =
    \frac{\sum_i (X_i-\Xbar)(Y_i-\bar Y)}
         {\sqrt{\sum_i(X_i-\Xbar)^2}\sqrt{\sum_i(Y_i-\bar Y)^2}}
\]
\begin{itemize}
    \item $-1 \le r \le 1$; unit-free and scale-invariant.
    \item $|r|=1$ $\Leftrightarrow$ the points lie exactly on a straight
          line.
    \item $r=0$: no linear relationship (nonlinear relationship still
          possible).
    \item $r$ is unchanged by any positive linear transform of $X$ or $Y$
          (e.g.\ changing units from cm to inches), and flips sign under a
          negative linear transform.
\end{itemize}
\note{Cauchy-Schwarz inequality is exactly why |r| <= 1: it states
|sum a_i b_i| <= sqrt(sum a_i^2) sqrt(sum b_i^2), applied here with
a_i = X_i - Xbar, b_i = Y_i - Ybar. Worth mentioning this as the underlying
mathematical reason, not just an empirical fact. Also stress correlation
does not imply causation - a classic caveat to state explicitly in any
exam answer discussing correlation.}
\end{frame}

% ---------------- SAMPLE COVARIANCE MATRIX ----------------
\subsection{Sample Covariance Matrix}

\begin{frame}{Sample Covariance Matrix}
For $p$-dimensional observations $\mathbf X_1,\dots,\mathbf X_n \in
\mathbb R^p$ with sample mean vector
$\bar{\mathbf X}=\frac1n\sum_i \mathbf X_i$:
\[
    \mathbf S = \frac{1}{n-1}\sum_{i=1}^n
    (\mathbf X_i-\bar{\mathbf X})(\mathbf X_i-\bar{\mathbf X})^\top
\]
In matrix form, with data matrix $\mathbf X$ ($n\times p$, rows =
observations) and centering matrix
$\mathbf H = \mathbf I_n - \frac1n \mathbf 1\mathbf 1^\top$:
\[
    \mathbf S = \frac{1}{n-1}\,\mathbf X^\top \mathbf H \mathbf X
\]
\begin{itemize}
    \item $\mathbf S$ is $p\times p$, \textbf{symmetric} and
          \textbf{positive semi-definite}.
    \item Diagonal entry $S_{jj}$ = sample variance of variable $j$;
          off-diagonal $S_{jk}$ = sample covariance of variables $j,k$.
\end{itemize}
\note{Positive semi-definiteness matters practically: for any vector a,
a^T S a = (1/(n-1)) sum_i [a^T(X_i - Xbar)]^2 >= 0, since it's a sum of
squares of the scalar projections a^T X_i. This guarantees any linear
combination of the variables has a non-negative estimated variance, as it
must. The centering matrix H is idempotent (H H = H) and symmetric,
which is why it appears naturally when you write out sums of squared
centered deviations in matrix form.}
\end{frame}

\begin{frame}{Sample Correlation Matrix}
Let $\mathbf D=\operatorname{diag}(\sqrt{S_{11}},\dots,\sqrt{S_{pp}})$.
\[
    \mathbf R = \mathbf D^{-1}\mathbf S\mathbf D^{-1}
\]
\begin{itemize}
    \item $\mathbf R$ has $1$'s on the diagonal, entries in $[-1,1]$
          elsewhere.
    \item Standardizes each variable, removing the effect of differing
          units/scales, so that variables can be compared directly.
    \item Both $\mathbf S$ and $\mathbf R$ are the starting point for
          multivariate techniques:
    \begin{itemize}
        \item \textbf{PCA}: eigen-decomposition of $\mathbf S$ (or
              $\mathbf R$) to find directions of maximal variance.
        \item \textbf{Mahalanobis distance}: uses $\mathbf S^{-1}$ to
              measure distance accounting for correlation between
              variables.
        \item Multivariate normal likelihood is parameterized directly by
              a covariance matrix.
    \end{itemize}
\end{itemize}
\note{PCA link is a good bridge to more advanced courses: the eigenvectors
of S give the principal component directions, and the eigenvalues give the
variance explained along each direction. Mention that S must be invertible
(full rank, requiring n > p typically) for Mahalanobis distance and many
multivariate procedures to be well-defined; when p is large relative to n,
S becomes singular or ill-conditioned, a known practical issue in
high-dimensional statistics.}
\end{frame}

% ============================================================
\section{Frequentist Statistics}
% ============================================================

\begin{frame}{Frequentist Statistics: Overview}
\Def{Under the frequentist interpretation, probability is the long-run
relative frequency of an event under (hypothetical) repeated sampling.
Population parameters ($\theta$) are treated as fixed, unknown constants;
randomness enters only through the sampling process, not through
$\theta$ itself.}
\begin{itemize}
    \item Contrast with Bayesian statistics, which treats $\theta$ itself as
          a random variable with a prior distribution.
    \item Core frequentist tools: point estimation, sampling distributions,
          bias/variance/MSE, consistency, confidence intervals, hypothesis
          testing.
    \item This module covers: sampling, MSE, consistency, confidence
          intervals, and parametric \& non-parametric estimation.
\end{itemize}
\note{It is worth explicitly contrasting frequentist and Bayesian
philosophies even briefly, since exam questions sometimes ask students to
articulate the difference: frequentist statements like a 95% CI are about
long-run coverage properties of the PROCEDURE across repeated samples, not
a probability statement about the specific fixed parameter given the one
observed data set. This distinction resurfaces explicitly in the
confidence interval section.}
\end{frame}

% ---------------- SAMPLING ----------------
\subsection{Sampling}

\begin{frame}{Sampling: Basic Concepts}
\begin{itemize}
    \item \textbf{Population}: entire set of units of interest.
          \textbf{Sample}: the observed subset.
    \item \textbf{Random (i.i.d.) sample}: $X_1,\dots,X_n$ independent,
          each distributed as the population distribution $F_\theta$.
    \item \textbf{Statistic}: a function $T(X_1,\dots,X_n)$ of the sample
          that does not depend on unknown parameters (e.g.\ $\Xbar$, $S^2$,
          sample max).
    \item \textbf{Sampling distribution}: the probability distribution of a
          statistic $T$, induced purely by the randomness in the sample.
    \item \textbf{Standard error}: the standard deviation of a statistic's
          sampling distribution, e.g.\ $\text{SE}(\Xbar)=\sigma/\sqrt n$.
\end{itemize}
\note{Distinguish clearly between a PARAMETER (fixed, unknown, a property
of the population, e.g. mu, sigma^2) and a STATISTIC (random, computable
from data, e.g. Xbar, S^2) and an ESTIMATE (the numeric value a statistic
takes for one particular observed data set). Confusing these three is one
of the most common sources of error in introductory exams.}
\end{frame}

\begin{frame}{Sampling Methods}
\begin{itemize}
    \item \textbf{Simple random sampling (SRS)}: every subset of size $n$
          is equally likely to be chosen.
    \item \textbf{Stratified sampling}: divide the population into
          homogeneous strata, sample independently within each — reduces
          variance when strata differ substantially.
    \item \textbf{Cluster sampling}: divide into clusters, randomly select
          whole clusters — cheaper logistically, but usually less efficient
          than SRS since units within a cluster tend to be similar.
    \item \textbf{Systematic sampling}: choose every $k$-th unit from an
          ordered list, with a random starting point.
    \item \textbf{With vs.\ without replacement}: sampling with replacement
          preserves exact independence; sampling without replacement from a
          finite population of size $N$ requires the \textbf{finite
          population correction} factor $\sqrt{(N-n)/(N-1)}$ multiplying the
          standard error.
\end{itemize}
\note{Finite population correction (FPC) approaches 1 (negligible) when
n << N, which is why in most textbook problems (infinite or very large
populations) it is simply ignored; it becomes important in survey sampling
of small, finite populations where you sample an appreciable fraction of
the whole population.}
\end{frame}

\begin{frame}{Sampling Distribution of the Mean}
If $X_1,\dots,X_n$ i.i.d.\ with mean $\mu$, variance $\sigma^2$:
\[
    \E[\Xbar]=\mu, \qquad \Var(\Xbar)=\frac{\sigma^2}{n}
\]
\begin{itemize}
    \item Population exactly Normal $\Rightarrow$
          $\Xbar \sim N(\mu,\sigma^2/n)$ \textbf{exactly}, for every $n$.
    \item Otherwise, by the CLT, for large $n$:
          $\Xbar \approx N(\mu,\sigma^2/n)$ approximately.
    \item \textbf{Rule of thumb}: CLT approximation is usually adequate for
          $n\gtrsim 30$ for moderately skewed populations; more skewed or
          heavy-tailed populations need larger $n$.
\end{itemize}
\note{Also worth mentioning Chebyshev's inequality here as a distribution-
free (non-normal-assuming) bound: P(|Xbar - mu| >= k * SE) <= 1/k^2. It is
much looser than the normal-based bound but requires no distributional
assumption beyond finite variance, illustrating the trade-off between
assumption strength and bound tightness that recurs throughout the
frequentist section.}
\end{frame}

% ---------------- MEAN SQUARE ERROR ----------------
\subsection{Mean Square Error}

\begin{frame}{Mean Square Error (MSE)}
For an estimator $\hat\theta$ of parameter $\theta$:
\[
    \MSE(\hat\theta)=\E\big[(\hat\theta-\theta)^2\big]
    = \Var(\hat\theta) + \big[\Bias(\hat\theta)\big]^2,
    \qquad \Bias(\hat\theta)=\E[\hat\theta]-\theta
\]
\textbf{Derivation:}
\begin{align*}
\E[(\hat\theta-\theta)^2]
&= \E\Big[\big((\hat\theta-\E[\hat\theta]) + (\E[\hat\theta]-\theta)\big)^2\Big]\\
&= \Var(\hat\theta) + 2(\E[\hat\theta]-\theta)\underbrace{\E[\hat\theta-\E[\hat\theta]]}_{=0}
   + \big(\E[\hat\theta]-\theta\big)^2\\
&= \Var(\hat\theta)+\Bias(\hat\theta)^2
\end{align*}
\note{This decomposition is arguably the single most important formula in
the whole module and reappears throughout applied statistics and machine
learning as the "bias-variance trade-off". Make sure students can both
recite it and reproduce the short derivation: it hinges only on adding and
subtracting E[theta-hat] inside the square and noting the cross term has
expectation zero because E[theta-hat - E[theta-hat]] = 0 by definition of
expectation.}
\end{frame}

\begin{frame}{MSE: Interpretation, Trade-offs, Example}
\begin{itemize}
    \item MSE is the standard criterion to compare competing estimators —
          smaller MSE is better.
    \item Unbiased estimator: $\Bias=0 \Rightarrow \MSE=\Var$.
    \item A slightly \textbf{biased} estimator can have a much smaller
          variance, giving \emph{lower total MSE} than an unbiased
          alternative — the bias--variance trade-off (e.g.\ ridge
          regression vs.\ ordinary least squares).
\end{itemize}
\Ex{Compare $\hat\sigma^2=\frac1n\sum(X_i-\Xbar)^2$ (biased, MLE) with
$S^2=\frac{1}{n-1}\sum(X_i-\Xbar)^2$ (unbiased). Under normality it can be
shown that $\hat\sigma^2$, despite its bias, has \emph{smaller} MSE than
$S^2$ for estimating $\sigma^2$ — biasedness alone does not make an
estimator worse overall.}
\note{This worked example is a favourite "trick" exam question: students
who only check unbiasedness will wrongly conclude S^2 is unambiguously
"better" than the MLE version; but MSE, the more complete criterion,
sometimes favours the biased estimator. Encourage computing both Var and
Bias^2 explicitly for full marks rather than asserting the conclusion.}
\end{frame}

% ---------------- CONSISTENCY ----------------
\subsection{Consistency}

\begin{frame}{Consistency of Estimators}
$\hat\theta_n$ is a \textbf{consistent} estimator of $\theta$ if it
converges in probability to $\theta$:
\[
    \hat\theta_n \xrightarrow{P} \theta
    \iff
    \forall \varepsilon>0,\quad \lim_{n\to\infty} P(|\hat\theta_n-\theta|>\varepsilon)=0
\]

\textbf{Sufficient condition (mean-square consistency):} if
$\Bias(\hat\theta_n)\to 0$ and $\Var(\hat\theta_n)\to 0$ as $n\to\infty$,
then $\MSE(\hat\theta_n)\to 0$, which implies $\hat\theta_n
\xrightarrow{P}\theta$ (by Chebyshev's inequality applied to
$\hat\theta_n - \theta$).

\begin{itemize}
    \item $\Xbar$ is consistent for $\mu$: $\Bias=0$,
          $\Var(\Xbar)=\sigma^2/n\to 0$.
    \item $S^2$ is consistent for $\sigma^2$.
    \item $\hat\sigma^2$ (MLE, divide by $n$) is also consistent for
          $\sigma^2$, even though biased for every finite $n$.
\end{itemize}
\note{Prove the sufficient condition explicitly if time permits: by
Chebyshev, P(|theta-hat - theta| > eps) <= MSE(theta-hat)/eps^2 -> 0 if
MSE -> 0. This directly links back to the MSE slide, reinforcing why MSE
convergence is a convenient, checkable route to establishing consistency
without working with the probability statement directly.}
\end{frame}

\begin{frame}{Consistency vs.\ Unbiasedness}
\begin{itemize}
    \item \textbf{Unbiasedness} is a \emph{finite-sample} property:
          $\E[\hat\theta_n]=\theta$ must hold for every fixed $n$.
    \item \textbf{Consistency} is an \emph{asymptotic} property: behavior
          only as $n\to\infty$.
    \item Possible combinations:
    \begin{itemize}
        \item Unbiased \emph{and} consistent — e.g.\ $\Xbar$, $S^2$.
        \item Biased but consistent — e.g.\ MLE variance
              $\hat\sigma^2=\frac1n\sum(X_i-\Xbar)^2$: biased for finite
              $n$, but bias $\to 0$ as $n\to\infty$.
        \item Unbiased but \emph{not} consistent — possible in
              pathological/contrived constructions (e.g.\ an estimator that
              ignores most of the data and has non-vanishing variance).
        \item Neither unbiased nor consistent — a poorly chosen estimator.
    \end{itemize}
\end{itemize}
\note{Give a concrete pathological example if asked: an estimator that
always just uses the FIRST observation, theta-hat = X_1, is unbiased for
mu (since E[X_1]=mu) but is NOT consistent, because its variance stays at
sigma^2 forever and never shrinks as n grows — it never "uses" the
additional data. This crisply separates the two concepts.}
\end{frame}

% ---------------- CONFIDENCE INTERVALS ----------------
\subsection{Confidence Intervals}

\begin{frame}{Confidence Intervals: Definition and Interpretation}
A $100(1-\alpha)\%$ confidence interval (CI) for $\theta$ is a random
interval $[L(\mathbf X), U(\mathbf X)]$ satisfying
\[
    P(L \le \theta \le U) = 1-\alpha
\]
\textbf{Correct frequentist interpretation:} if the sampling procedure were
repeated many times and a CI constructed each time, approximately
$100(1-\alpha)\%$ of those intervals would contain the true, fixed
$\theta$. The interval is random; $\theta$ is fixed.

\textbf{Common wrong interpretation to avoid:} ``there is a 95\% probability
that $\theta$ lies in \emph{this particular} computed interval'' — once
data are observed and $L,U$ are specific numbers, $\theta$ either is or is
not in $[L,U]$; there is no probability left in a frequentist framework.
\note{This slide addresses the single most commonly mis-stated concept in
introductory statistics. If an exam asks you to interpret a specific
computed 95% CI, always phrase it in terms of the long-run performance of
the PROCEDURE ("we are 95% confident" as a statement about the method),
never as a probability statement about the fixed parameter given this one
realized interval.}
\end{frame}

\begin{frame}{CI for the Mean: $\sigma$ Known}
If $X_1,\dots,X_n\stackrel{iid}\sim N(\mu,\sigma^2)$ (or $n$ large via CLT)
and $\sigma$ is \textbf{known}:
\[
    \Xbar \pm z_{\alpha/2}\,\frac{\sigma}{\sqrt n}
\]
where $z_{\alpha/2}$ is the upper-$\alpha/2$ quantile of $N(0,1)$
($z_{0.025}=1.96$ for a 95\% CI; $z_{0.005}=2.576$ for 99\%).

\textbf{Derivation:} $Z=\dfrac{\Xbar-\mu}{\sigma/\sqrt n}\sim N(0,1)$
(this $Z$ is a \textbf{pivotal quantity} — see later slide). Then
$P(-z_{\alpha/2}\le Z\le z_{\alpha/2})=1-\alpha$; rearranging the inequality
to isolate $\mu$ gives the interval above.
\note{Walk through the algebra of isolating mu from the inequality
explicitly in class: -z <= (Xbar - mu)/(sigma/sqrt n) <= z rearranges to
Xbar - z*sigma/sqrt(n) <= mu <= Xbar + z*sigma/sqrt(n). This exact
rearrangement technique (start from a known probability statement about a
pivot, then isolate theta) is the GENERAL METHOD used for every CI in this
module, not just this one case.}
\end{frame}

\begin{frame}{CI for the Mean: $\sigma$ Unknown (t-interval)}
Replace $\sigma$ with sample standard deviation $S$; use Student's
$t$-distribution with $n-1$ degrees of freedom:
\[
    \Xbar \pm t_{\alpha/2,\,n-1}\,\frac{S}{\sqrt n}
\]
\begin{itemize}
    \item $T=\dfrac{\Xbar-\mu}{S/\sqrt n}\sim t_{n-1}$ when the population
          is normal.
    \item $t$-distribution has heavier tails than $N(0,1)$
          $\Rightarrow t_{\alpha/2,n-1} > z_{\alpha/2}$, reflecting the
          extra uncertainty from also having to estimate $\sigma$.
    \item As $n\to\infty$, $t_{n-1}\to N(0,1)$, so the $t$- and
          $z$-intervals converge for large samples.
\end{itemize}
\note{Emphasize WHY the t-distribution is needed at all: once you replace
the known sigma with the estimated S, the ratio (Xbar-mu)/(S/sqrt n) picks
up extra variability from S itself being random, and this extra
variability is exactly what gives the t-distribution its heavier tails
relative to the normal. Degrees of freedom n-1 again reflects that one
degree of freedom was "used up" estimating mu via Xbar inside S^2 itself.}
\end{frame}

\begin{frame}{CI for a Proportion}
For large $n$, with sample proportion $\hat p = X/n$ ($X$ = number of
successes out of $n$ Bernoulli trials):
\[
    \hat p \pm z_{\alpha/2}\sqrt{\frac{\hat p(1-\hat p)}{n}}
\]
based on the normal approximation to the Binomial distribution, generally
considered valid when $n\hat p \ge 5$ and $n(1-\hat p)\ge 5$.

\Note{An alternative, often more accurate for small/moderate $n$, is the
\textbf{Wilson score interval}, which does not rely on the same crude
normal approximation and performs better near $\hat p$ close to 0 or 1.}
\note{Point out the naive Wald interval above can behave badly (even
producing bounds outside [0,1]) when p-hat is near 0 or 1 or n is small —
worth mentioning the Wilson interval by name as a superior alternative
even if its formula is not required to be memorized, since it shows
awareness of the practical limitations of the basic method.}
\end{frame}

\begin{frame}{CI for Variance}
If $X_1,\dots,X_n\stackrel{iid}\sim N(\mu,\sigma^2)$:
\[
    \frac{(n-1)S^2}{\sigma^2}\sim \chi^2_{n-1}
\]
$100(1-\alpha)\%$ CI for $\sigma^2$:
\[
    \left(\frac{(n-1)S^2}{\chi^2_{\alpha/2,\,n-1}},\ \
    \frac{(n-1)S^2}{\chi^2_{1-\alpha/2,\,n-1}}\right)
\]
where $\chi^2_{\alpha/2,n-1}$ is the upper-$\alpha/2$ quantile of the
chi-square distribution with $n-1$ degrees of freedom.
\note{Note the asymmetry: unlike the mean's CI, this interval is NOT of
the simple "estimate plus-or-minus margin" form, because the chi-square
distribution is itself asymmetric (skewed right), and dividing by a larger
chi-square quantile gives the SMALLER endpoint while dividing by the
smaller quantile gives the LARGER endpoint — students often flip these by
mistake, so double check which chi-square quantile goes with which bound.}
\end{frame}

\begin{frame}{General Construction: Pivotal Quantities}
\Def{A pivotal quantity is a function $Q(\mathbf X,\theta)$ of the data and
the parameter whose sampling distribution does not depend on $\theta$ or
any other unknown parameter.}
\begin{itemize}
    \item \textbf{Recipe:} (1) find a pivot $Q$; (2) using its known
          distribution, find $a,b$ with $P(a\le Q\le b)=1-\alpha$;
          (3) algebraically invert the inequality to isolate $\theta$.
    \item Pivots seen so far: $Z=\frac{\Xbar-\mu}{\sigma/\sqrt n}$,
          $T=\frac{\Xbar-\mu}{S/\sqrt n}$, $\chi^2=\frac{(n-1)S^2}{\sigma^2}$.
    \item CI width typically shrinks like $1/\sqrt n$ — more data gives a
          more precise (narrower) interval, all else equal.
\end{itemize}
\note{The pivotal-quantity method is the unifying, general-purpose
technique behind every specific CI formula on the previous slides — if a
student forgets a formula in an exam, they can, in principle, re-derive it
from scratch given (a) an appropriate pivot with a KNOWN distribution not
involving the unknown parameter, and (b) tables/quantiles of that
distribution. This is worth stressing as the "big picture" takeaway of the
whole confidence-interval section.}
\end{frame}

% ---------------- PARAMETRIC ESTIMATION ----------------
\subsection{Parametric Model Estimation}

\begin{frame}{Parametric Estimation: Setup}
\Def{A parametric model assumes data come from a distribution family
$\{f_\theta:\theta\in\Theta\}$ indexed by a finite-dimensional parameter
$\theta\in\Theta\subseteq\mathbb R^k$ (e.g.\ Normal$(\mu,\sigma^2)$,
Binomial$(n,p)$, Poisson$(\lambda)$, Exponential$(\lambda)$).}
Main estimation methods covered:
\begin{enumerate}
    \item Method of Moments (MoM)
    \item Maximum Likelihood Estimation (MLE)
    \item (Brief mention) Least Squares — for regression-type models
    \item (Brief mention) Bayesian estimation — for contrast; not strictly
          frequentist since it treats $\theta$ as random with a prior.
\end{enumerate}
\note{Reiterate the key defining feature of "parametric": the entire
unknown distribution is captured by a FINITE, fixed-length list of numbers
(theta), so estimation reduces to estimating that finite list, in contrast
to the non-parametric setting later where the "parameter" is effectively
an entire unknown function.}
\end{frame}

\begin{frame}{Method of Moments}
\textbf{Idea:} equate population moments $\mu_k=\E[X^k]$ to sample moments
$m_k=\frac1n\sum_i X_i^k$ and solve for $\theta$.

\Ex{Normal$(\mu,\sigma^2)$:}
\[
    \hat\mu_{MoM}=\Xbar, \qquad
    \hat\sigma^2_{MoM}=\frac1n\sum_i (X_i-\Xbar)^2
\]
\Ex{Exponential$(\lambda)$, where $\E[X]=1/\lambda$:}
\[
    \hat\lambda_{MoM}=1/\Xbar
\]
\Ex{Uniform$(0,\theta)$, where $\E[X]=\theta/2$:}
\[
    \hat\theta_{MoM}=2\Xbar
\]
\Note{MoM estimators are simple to compute but can sometimes fall outside
the valid parameter space (e.g.\ $2\Xbar$ could in principle be less than
$X_{(n)}$'s observed maximum) and are generally less statistically
efficient than MLE.}
\note{The Uniform(0,theta) example is a nice one to contrast with MLE
later: the MLE for this model turns out to be the sample maximum X_(n),
which is quite different in character from the MoM estimator 2*Xbar — a
good illustration that different estimation principles can give genuinely
different-looking estimators for the same parameter.}
\end{frame}

\begin{frame}{Maximum Likelihood Estimation (MLE)}
\textbf{Likelihood function:}
\[
    L(\theta\mid x_1,\dots,x_n)=\prod_{i=1}^n f_\theta(x_i)
\]
\textbf{MLE:}
\[
    \hat\theta_{MLE}=\arg\max_{\theta\in\Theta} L(\theta)
    =\arg\max_\theta \ell(\theta), \qquad \ell(\theta)=\log L(\theta)
\]
found (when $\ell$ is differentiable) by solving the \textbf{score
equation} $\partial\ell(\theta)/\partial\theta=0$.

\textbf{Properties (under regularity conditions), as $n\to\infty$:}
\begin{itemize}
    \item Consistent: $\hat\theta_{MLE}\xrightarrow{P}\theta$.
    \item Asymptotically normal and asymptotically efficient (attains the
          Cram\'er--Rao lower bound in the limit).
    \item \textbf{Invariance property}: the MLE of $g(\theta)$ is
          $g(\hat\theta_{MLE})$ for any function $g$.
\end{itemize}
\note{The invariance property is powerful and easy to forget: e.g. if you
have the MLE for sigma^2, the MLE for sigma (the standard deviation) is
simply the square root of the MLE for sigma^2 - you do NOT need to
re-derive a separate likelihood maximization for every transformed
parameter of interest.}
\end{frame}

\begin{frame}{MLE Example: Normal Distribution}
For $X_1,\dots,X_n\stackrel{iid}\sim N(\mu,\sigma^2)$:
\[
    \ell(\mu,\sigma^2)=-\frac n2\log(2\pi\sigma^2)
    -\frac{1}{2\sigma^2}\sum_{i=1}^n (x_i-\mu)^2
\]
Setting $\partial \ell/\partial\mu=0$ and $\partial\ell/\partial\sigma^2=0$
and solving jointly:
\[
    \hat\mu_{MLE}=\Xbar, \qquad
    \hat\sigma^2_{MLE}=\frac1n\sum_{i=1}^n (x_i-\Xbar)^2
\]
Note $\hat\sigma^2_{MLE}$ is \textbf{biased} (divides by $n$, not $n-1$)
but is still \textbf{consistent}.
\end{frame}

\begin{frame}{MLE Example: Bernoulli/Binomial and Poisson}
\Ex{Bernoulli$(p)$, data $X_1,\dots,X_n\in\{0,1\}$:}
\[
    \ell(p)=\left(\sum_i x_i\right)\log p + \left(n-\sum_i x_i\right)\log(1-p)
    \ \Rightarrow\ \hat p_{MLE}=\frac1n\sum_i X_i = \Xbar
\]
\Ex{Poisson$(\lambda)$:}
\[
    \ell(\lambda)= -n\lambda + \Big(\sum_i x_i\Big)\log\lambda - \sum_i\log(x_i!)
    \ \Rightarrow\ \hat\lambda_{MLE}=\Xbar
\]
\begin{itemize}
    \item In both cases the MLE coincides exactly with the sample mean and
          with the Method-of-Moments estimator — this is common (though not
          universal) for simple one-parameter exponential-family models.
\end{itemize}
\note{Deliberately chosen so MLE = MoM = sample mean in both cases -
useful for building intuition, but stress explicitly this coincidence does
NOT hold in general (e.g. Uniform(0,theta) from two slides ago, where MLE
= X_(n) but MoM = 2*Xbar - very different answers).}
\end{frame}

\begin{frame}{Fisher Information and the Cram\'er--Rao Lower Bound}
\textbf{Fisher information} (single observation):
\[
    I(\theta)=\E\left[\left(\frac{\partial\log f_\theta(X)}{\partial\theta}\right)^2\right]
    = -\E\left[\frac{\partial^2\log f_\theta(X)}{\partial\theta^2}\right]
\]
\textbf{Cram\'er--Rao lower bound (CRLB):} for any unbiased estimator
$\hat\theta$ based on $n$ i.i.d.\ observations,
\[
    \Var(\hat\theta)\ \ge\ \frac{1}{n\,I(\theta)}
\]
An estimator achieving this bound with equality is called
\textbf{efficient}. MLE is asymptotically efficient (achieves the bound as
$n\to\infty$, under regularity conditions).
\note{Two equivalent formulas for Fisher information are given (expected
square of the score, and negative expected second derivative /
"curvature" of the log-likelihood) - the second is usually algebraically
easier to compute in practice for standard distributions. CRLB gives an
absolute floor on how small the variance of ANY unbiased estimator can be;
this is why "achieving the CRLB" is treated as the gold standard of
estimator efficiency, analogous to how MSE was the gold standard
comparison criterion earlier.}
\end{frame}

% ---------------- NON-PARAMETRIC ESTIMATION ----------------
\subsection{Non-parametric Model Estimation}

\begin{frame}{Non-parametric Estimation: Setup}
\Def{A non-parametric model makes no finite-dimensional assumption about
the form of $f$ or $F$; the object being estimated (the whole function) is
effectively infinite-dimensional.}
\textbf{Trade-off vs.\ parametric methods:}
\begin{itemize}
    \item \textbf{Parametric}: highly efficient \emph{if} the assumed model
          is correct; can be badly biased if the model is wrong
          (misspecified).
    \item \textbf{Non-parametric}: robust to misspecification (fewer
          assumptions to get wrong); typically requires more data and
          converges more slowly than a correctly specified parametric
          model would.
\end{itemize}
\note{This slide's trade-off framing directly parallels the earlier
bias-variance framing of MSE: choosing a parametric model is like
accepting some risk of bias (if wrong) in exchange for lower variance
(fewer things to estimate); choosing non-parametric methods reduces
model-misspecification bias risk at the cost of needing more data /
having higher variance for a given sample size.}
\end{frame}

\begin{frame}{Empirical Distribution Function}
\[
    \hat F_n(x)=\frac1n\sum_{i=1}^n \mathbf 1\{X_i\le x\}
\]
\begin{itemize}
    \item Natural non-parametric estimator of the true cdf $F$: a step
          function jumping by $1/n$ at each observed data point.
    \item For fixed $x$: $n\hat F_n(x)\sim \text{Binomial}(n,F(x))$, so
          $\hat F_n(x)$ is \textbf{unbiased} with
          $\Var(\hat F_n(x))=F(x)(1-F(x))/n$.
    \item \textbf{Glivenko--Cantelli theorem} (``fundamental theorem of
          statistics''):
    \[
        \sup_x |\hat F_n(x)-F(x)| \xrightarrow{a.s.} 0 \quad\text{as } n\to\infty
    \]
    i.e.\ uniform, almost-sure convergence of the entire estimated curve to
    the true cdf, not just pointwise convergence.
\end{itemize}
\note{Glivenko-Cantelli is stronger than simple pointwise consistency at
each x (which follows immediately from the LLN applied to the indicator
variables) - it guarantees the WORST-CASE gap across all x simultaneously
vanishes, which is what justifies using the empirical cdf as a reliable
stand-in for the true cdf uniformly, e.g. in goodness-of-fit tests like
the Kolmogorov-Smirnov test, which is built directly on this supremum
statistic.}
\end{frame}

\begin{frame}{Kernel Density Estimation (KDE)}
\[
    \hat f_h(x)=\frac{1}{nh}\sum_{i=1}^n K\!\left(\frac{x-X_i}{h}\right)
\]
\begin{itemize}
    \item $K(\cdot)$: kernel function (e.g.\ Gaussian, Epanechnikov),
          typically symmetric and integrating to 1.
    \item $h>0$: \textbf{bandwidth}, controlling smoothness — the
          continuous analogue of histogram bin width.
    \item \textbf{Bias--variance trade-off in $h$}: small $h$ gives low bias
          but high variance (wiggly estimate); large $h$ gives high bias
          (oversmoothed) but low variance.
    \item KDE generalizes the histogram by replacing the rigid, hard-edged
          bin-membership indicator with a smooth, weighted kernel — this
          removes the artificial discontinuities (jumps) present in a
          histogram.
\end{itemize}
\note{Optimal bandwidth selection (e.g. via cross-validation, or rules
like Silverman's rule of thumb h = 1.06 * min(s, IQR/1.34) * n^{-1/5}) is
an important practical topic; the exponent -1/5 for KDE bandwidth
(compared to -1/3 for histogram bin width in Scott's rule) reflects that
KDE's smoother construction allows a different, in fact better,
asymptotic convergence rate for the estimated density.}
\end{frame}

\begin{frame}{Non-parametric Regression}
\textbf{Nadaraya--Watson kernel regression estimator:}
\[
    \hat m(x)=\frac{\sum_{i=1}^n K\!\left(\frac{x-X_i}{h}\right) Y_i}
                    {\sum_{i=1}^n K\!\left(\frac{x-X_i}{h}\right)}
\]
\begin{itemize}
    \item A kernel-weighted local average of the $Y_i$ near $x$, estimating
          $m(x)=\E[Y\mid X=x]$ without assuming any specific functional
          form (e.g.\ not assuming linearity as ordinary linear regression
          does).
    \item Related non-parametric regression methods: local polynomial
          regression, smoothing splines, $k$-nearest-neighbours regression.
    \item \textbf{Convergence rates}: non-parametric estimators typically
          converge more slowly than correctly specified parametric
          estimators — e.g.\ KDE/Nadaraya--Watson rate is roughly
          $n^{-2/5}$, versus $n^{-1/2}$ for a correctly specified parametric
          MLE — the statistical price paid for not assuming a model form.
\end{itemize}
\note{This n^{-2/5} vs n^{-1/2} rate comparison is the clean, quantitative
punchline of the entire "parametric vs non-parametric" discussion running
through this whole section: parametric models, when correct, converge
faster (the classical "root-n" rate), while non-parametric methods
converge more slowly in exchange for making far weaker assumptions about
the underlying model - a foundational idea that carries directly into
modern non-parametric and machine learning methods.}
\end{frame}

\begin{frame}{Bootstrap: A Non-parametric Resampling Method}
\begin{itemize}
    \item \textbf{Idea}: treat the observed empirical distribution
          $\hat F_n$ as a stand-in for the true (unknown) $F$, and simulate
          the sampling distribution of a statistic by repeatedly resampling
          \emph{with replacement} from the observed data.
    \item \textbf{Procedure}: draw $B$ resamples of size $n$ from
          $X_1,\dots,X_n$ (with replacement); compute the statistic of
          interest $\hat\theta^{*}_b$ on each resample; use the empirical
          distribution of $\{\hat\theta^*_1,\dots,\hat\theta^*_B\}$ to
          approximate the sampling distribution of $\hat\theta$, estimate
          its standard error, or build a confidence interval.
    \item Useful when the sampling distribution of a statistic has no
          simple closed-form expression (e.g.\ for the sample median or a
          complicated function of parameters).
\end{itemize}
\note{The bootstrap is a fully non-parametric, computational alternative
to the analytic pivotal-quantity method used earlier for constructing
confidence intervals - instead of relying on a KNOWN distribution for a
pivot (like Z, T, or chi-square), the bootstrap uses the computer to
approximate the sampling distribution directly from resampled data. This
makes a nice closing bridge between the classical CI section and modern
computational statistics.}
\end{frame}

% ============================================================
\section{Summary}
% ============================================================

\begin{frame}{Module Summary}
\begin{columns}[T]
\column{0.5\textwidth}
\textbf{Descriptive Statistics}
\begin{itemize}
    \item Histogram (bin-width rules, density estimate)
    \item Sample mean $\Xbar$, sample variance $S^2$
    \item Skewness and kurtosis (shape)
    \item Order statistics, quantiles, IQR, boxplot
    \item Sample covariance, correlation
    \item Sample covariance/correlation matrix, link to PCA
\end{itemize}
\column{0.5\textwidth}
\textbf{Frequentist Statistics}
\begin{itemize}
    \item Sampling, sampling distributions, standard error
    \item MSE = Variance + Bias$^2$
    \item Consistency vs.\ unbiasedness
    \item Confidence intervals via pivotal quantities
    \item Parametric estimation: MoM, MLE, Fisher info, CRLB
    \item Non-parametric: EDF, Glivenko--Cantelli, KDE,
          Nadaraya--Watson, bootstrap
\end{itemize}
\end{columns}
\note{Use this slide as a final review checklist: for a thorough exam
preparation pass, a student should be able to state a precise definition,
write the key formula, and give at least one derivation or worked example
for every single bullet point listed here, including the newly expanded
skewness and kurtosis subsection.}
\end{frame}

\begin{frame}
\centering
\Huge{Thank You}\\[10pt]
\end{frame}

\end{document}